\section{Введение}
\label{sec:Chapter0}

\subsection*{Актуальность}

Сердечно-сосудистые заболевания~--- ведущая причина смертности в~мире:
по~данным ВОЗ, они ежегодно уносят более~17~миллионов жизней~\cite{who2021}.
Среди них особое место занимают нарушения ритма~--- аритмии~--- которые нередко
протекают без~явных симптомов, однако при~несвоевременной диагностике могут
приводить к~инфаркту миокарда, инсульту или~внезапной сердечной смерти.
Основным методом их~выявления служит электрокардиограмма~(ЭКГ).

При~суточном холтеровском мониторировании за~24~часа регистрируется до~120\,000
сердечных сокращений.
Ручная обработка таких записей~--- трудоёмкая задача, а~усталость специалиста
неизбежно влечёт ошибки.
Автоматическая классификация ЭКГ-сигналов с~помощью алгоритмов машинного обучения
позволяет снять эту нагрузку и~ускорить диагностику.

Дополнительную сложность создаёт \textit{нестационарность} ЭКГ: амплитуда,
форма и~частота сигнала меняются во~времени под~влиянием физической нагрузки,
дыхания и~эмоционального состояния пациента.
Для~анализа таких сигналов эффективен вейвлет-анализ~\cite{zakharova2018}~---
инструмент частотно-временного описания, позволяющий видеть
и~частотный состав, и~момент возникновения каждой составляющей.
За~последние годы глубокое обучение существенно продвинулось в~этой задаче
\cite{hannun2019}, хотя классические ансамблевые методы~--- прежде всего
случайный лес~--- по-прежнему ценны своей интерпретируемостью и~стойкостью к~шуму.

\subsection*{Цель и задачи работы}

Цель настоящей работы~--- исследование и~сравнение методов автоматической
классификации кардиоциклов на~базе данных MIT-BIH Arrhythmia Database, включая
оценку устойчивости моделей к~реалистичным артефактам записи.

Для~достижения этой цели в~работе решаются следующие задачи.
Разрабатывается конвейер предобработки данных~--- извлечение кардиоциклов,
нормализация и~стратифицированное разбиение при~наличии выраженного дисбаланса
классов.
Реализуются и~обучаются три модели: случайный лес, одномерная сверточная
нейронная сеть и~остаточная нейронная сеть; для~борьбы с~дисбалансом применяется
взвешенная функция потерь.
Проводится частотно-временной вейвлет-анализ сигналов и~проверяется гипотеза
об~улучшении качества классификации при~добавлении вейвлет-признаков.
Оценивается устойчивость моделей к~гауссову шуму, дрейфу изолинии, сетевой помехе
50~Гц и~мышечному шуму.
Достоверность результатов подтверждается пятикратной стратифицированной
кросс-валидацией.

\subsection*{Структура работы}

Работа состоит из~пяти разделов.
В~разделе~\ref{sec:Chapter1} формализуется постановка задачи и~подробно описываются
данные и~конвейер предобработки.
Раздел~\ref{sec:Chapter2} посвящен теоретическому описанию использованных моделей
и~методов вейвлет-анализа.
В~разделе~\ref{sec:Chapter3} приводятся результаты классификации на~чистых данных~---
матрицы ошибок, кривые обучения и~важность признаков.
Раздел~\ref{sec:Chapter4} содержит анализ устойчивости моделей к~четырём типам
помех и~результаты кросс-валидации.
В~заключении~(\ref{sec:Chapter5}) формулируются основные выводы и~намечаются
направления дальнейших исследований.
